\documentclass[10pt]{article}

% graphicx package, useful for including eps, pdf, and raster graphics
\usepackage{graphicx}
\DeclareGraphicsExtensions{.png,.png,.jpg,.pdf}

% basic packages
\usepackage{fancyhdr}
\usepackage{color}
\usepackage{parskip}
\usepackage{float}
\usepackage{microtype}
\usepackage{url}
\usepackage{hyperref}
\usepackage{booktabs}
\usepackage{makecell}
\usepackage{multirow}
\usepackage{pbox}
\usepackage{enumitem}

% page layout
\usepackage[top=1in, bottom=1in, left=0.75in, right=0.75in, includefoot, heightrounded]{geometry}

% Useful for seeing the layout as you tweak it.
%\usepackage{showframe}

\setlength{\headheight}{1.1in}
\pagestyle{fancy}
\fancyhf{}

% fonts
\usepackage{fontspec}

\usepackage{polyglossia}
\setmainlanguage{amharic}
\setotherlanguage{english}
\newfontfamily\ethiopicfont[
  Path        = /src/fonts/,
  Extension   = .ttf,
  UprightFont = jiret,
  BoldFont    = washrab ]{}


\setmainfont{DroidSans}[
  Extension   = .ttf,
  UprightFont = *,
  BoldFont    = *-Bold ]
\setmonofont{DroidSansMono}[
  Extension   = .ttf,
  UprightFont = * ]


\hypersetup{colorlinks,urlcolor=blue}
\urlstyle{same}

\lhead{\includegraphics[height=1in,clip]{logo}}
\rhead{\pbox[b][1in][c]{\textwidth}{
    \small
    \textbf{መርመራ ዘካየደ ላባራቶሪ እዚ ዝስዕብ’ዩ:}\\
    Northwest Genomics Center\\
    University of Washington School of Medicine\\
    Genome Sciences Box 355065\\
    3720 15th Ave NE\\
    Seattle, WA 98195-5065}}

\newcommand{\PageLine}{\rule{\textwidth}{0.25mm}}

\begin{document}

\begin{description}[font=\normalfont,align=left,labelwidth=10em]
\item [\textbf{\VAR{pat_name|e}}]
\item [\textbf{ዕለተ ልደት:}] \VAR{birth_date|e}
\item [\textbf{መርኣያ:}]
  %- if swab_type == "ans"
  XXXXXXXX TODO
  %- elif swab_type == "mtb"
  ፍርቀ-ኣፍንጫ፡
  %- else
  \errmessage{Invalid swab_type "\VAR{swab_type}"}
  %- endif

  self-collected

\item [\textbf{ናይ ምውሳድ ባርኮድ:}] \VAR{qrcode|e}
\item [\textbf{ናይ ምውሳድ ዕለት:}] \VAR{collect_ts|e}
\item [\textbf{ናይ ምሕባር ዕለት:}]
  %- if status_code == "never-tested"
  ተመርሚረ ኣይፈልጥን
  %- else
  \VAR{result_ts|e}
  %- endif
\end{description}

\PageLine

\begin{center}
\Large
\textbf{ውጽኢት ምርምር – ነዚ ውጽኢት'ዚ ንወሃቢ ሕክምናዊ ኣገልግሎትኩም ክተካፍልዎ ትተባብዑ።}
\end{center}

\bigskip

\large \underline{\textbf{ብዛዕባ እዚ መርመራ'ዚ}}

\begin{itemize}
\item

  ኮቪድ-19 SARS-CoV-2 \textbf{ተባሂሉ ብዝጽዋዕ} ቫይረስ ዝስዕብ \textbf{ሕማም} እዩ። እዚ ናይ
  መጽናዕቲ መርመራ'ዚ ነቲ ቫይረስ ኣብቲ ካብ ኣፍንጫ ዝተወስደ መርኣያ ኮቪድ-19 እንተልይኩም ዝምርምር እዩ።
  እቲ ተክኒካዊ ስያመ ነዚ ዓይነት መርመራ'ዚ “reverse-transcription polymerase chain reaction
  (ግልባጥ-ቅዳሕ ፖሊመረይስ ናይ ሰንሰለታዊ መልሰ-ትግባር)”ወይ ብሓጺር “RT-PCR” እዩ።

\end{itemize}

\bigskip

\large \underline{\textbf{ኣገደስቲ መዘኻኸሪታት}}

ቅድሚ ንውጽኢት ናይ መጽናዕቲ መርመራኹም ምርኣይኩም፡ ነዚ ዝስዕብ ክትርእይዎ ኣለኩም:

\begin{itemize}
\item

  እቲ \textbf{ዝወሰድክሞ ዓይነት} መርኣያ ከምኡ ድማ እቲ \textbf{ናይ መርመራ ላባራቶሪ} ኣካል
  \textbf{መጽናዕቲ መርመራ እዮም} እዚ ጸብጻብ'ዚ ከኣ ንውጽኢት ናይ መጽናዕቲ
  \textbf{መርመራኹም ዝገልጽ'ዩ።}

\item
  %- if staff_observed == "f"
  ነዚ መርኣያ'ዚ ወሲድኩምን ብገዛእ ርእስኹም (ወይ ንኻልእ ሰብ) ምልክት ገይርክምሉን ከምኡ ድማ ናብቲ ላባራቶሪ
  ብመገዲ ናይ ምብጽጻሕ ኣገልግሎት መሊስክሞ። ብበዓል ሞያ ጥዕና ከም ነርስ ወይ ዶክተር ዝተወስደ፡ ምልክት
  ዝተገብረሉ፡ ወይ ዝተገዓዓዘ ኣይኮነን።

  %- elif staff_observed == "t"
  ነዚ መርኣያ'ዚ ባዕልኻ (ወይ ንኻልእ ሰብ) ወሲድካዮ ምልክት ተገይርሉን ናብቲ ላባራቶሪ ከኣ ብመገዲ መልእኽተኛ
  ወይ ካልእ ዓይነት በዓል ሞያ ጥዕና ዘይኮነ ሰብ ተመሊሱን። ብከም ዶክተር ወይ ነርስ ዝኣመሰለ ወረቐት ስራሕ ዘለዎ
  በዓል ሞያ ኣይተወስደን ወይ ኣይተገዓዓዘን። 

  %- else
  \errmessage{Unknown value: staff_observed = "\VAR{staff_observed}"}

  %- endif

  እዚ ንውጽኢት መርመራኹም ክጸሉ ይኽእል፡ ወሃቢ ሕክምናዊ ኣገልግሎት ደጊሙ ክምርምረኩም ክደሊ ከኣ ይኽእል።


\item

  ኣካል ናይዚ ናይ መጽናዕቲ \textbf{መርመራ'ዚ ምዃን} ንንቡር ሕክምናዊ ክንክንኩም ዝትክእ ኣይኮነን። እንተድኣ
  ኣዚኹም ሓሚምኩም ኣለኹም ወይ ብዛዕባ ውጽኢት መርመራኹም ስክፍታ እንተልይኩም፡ ንዶክተር ወይ ወሃቢ
  ሕክምናዊ ኣገልግሎት ክትረኽብዎ ኣለኩም።

\end{itemize}

\bigskip
\PageLine

\large \underline{\textbf{ናይ SARS-CoV-2 (COVID-19) ናይ መጽናዕቲ ውጽኢት መርመራኹም}}

%- if status_code == "negative"
ሕማም ኣይተረኽበን (ኣሉታዊ)\\
%- elif status_code == "positive"
ሕማም ተረኺቡ (ኣወንታዊ)\\
%- elif status_code == "inconclusive"
ናብ መደምደምታ ክብጻሕ ኣይተኻእለን\\
%- elif status_code == "never-tested"
ክምርመር ኣይተኻእለን\\
%- else
\VAR{0/0} %# poor man's raise
%- endif

\PageLine
\bigskip

\large \underline{\textbf{ምርዳእ ውጽኢትኩም}}

እዚ ሰሌዳ'ዚ ንኽሎም መርመራኹም ከርእዮም ዝኽእል ውጽኢታት ይህብ ነፍስ ወከፍ እዞም ውጽኢታት እዚኦም እንታይ
ማለት ምዃኖም ከኣ ይገልጽ። ነዚ ናይ መጽናዕቲ ውጽኢት መርመራ'ዚ \textbf{ንዶክተርኩም ወይ ወሃቢ ሕክምናዊ
ኣገልግሎትኩም} ንምርኣይ ክተሕትምዎ ወይ ኢመይል ክትገብርዎ ትኽእሉ።

ናይ መርመራኹም ውጽኢት ብዘየገድስ፡ እንተድኣ ሕማም ይስምዓኩም ኣሎ ኾይኑ ኣብ ገዛ ክትኮኑን ኣብ ጥቓ ካልኦት
ሰባት ምዃን ክተወግዱን ኣለኩም። \textbf{እንተድኣ ምልክታት ሕማምኩም ቀጺሎም፡ ኣዚኹም ሓሚምኩም
እንተሊኹም ወይ ብዛዕባ ውጽኢት መርመራኹም ስክፍታ እንተልይኩም፡ ንዶክተር ወይ ወሃቢ ሕክምናዊ ኣገልግሎት
ክትረኽብዎ ኣለኩም።}

\begin{itemize}


\item

  \textbf{እንተድኣ መርመራኹም “ኣይተረኽበን (ኣሉታዊ)”} ዝብል ኮይኑ፡ እዚ ማለት ላባራቶሪና
  ነቲ ኮቪድ-19 ዘስዕብ ቫይረስ ኣብቲ ዝሰደድኩምልና መርኣያ ኣይረከበን ማለት’ዩ። ፍጹም ዝብሃል መርመራ የለን፡
  ዋላ’ውን መርመራኹም ኣሉታዊ ክንሱ እቲ ቫይረስ ክህልወኩም ንእሽቶ ተኽእሎ ኣሎ።

\item

  \textbf{እንተድኣ መርመራኹም “ተረኺቡ (ኣወንታዊ)”} ዝብል ኮይኑ፡ እዚ ማለት ላባራቶሪና
  ነቲ ኮቪድ-19 ዘስዕብ ቫይረስ ኣብቲ ዝሰደድኩምልና መርኣያ ረኺቡ ማለት’ዩ። ኣብ ገዛኹም ክትኮኑን ኣብ ጥቓ
  ካልኦት ሰባት ምዃን ክተወግዱን ኣለኩም፡ ብፍላይ እንተድኣ ሕማም ይስምዓኩም ኣሎ ኾይኑ።

  ውጽኢትኩም ነቲ ክፍለ-ግዝኣት ዋሽንግተን ከምኡ ድማ ንከባብያዊ ክፍልታት ጥዕና ክካፈል እዩ። ነዞም ውጽኢታት
  እዚኦም ባዕልኹም ምሕባር  \textbf{ኣየድልየኩምን} እዩ። እቲ ክፍሊ ጥዕና ብዛዕባ ሕማምኩምን ኣብ
  ሓሚምክምሉ ዝነበርኩም እዋን ኣብ ጥቓኹም ብዛዕባ ዝነበሩ ሰባትን ሕቶታት ንምሕታት ክድውለልኩም ይኽእል።

\item

  \textbf{እንተድኣ እቲ መርመራ “ናብ መደምደምታ ክብጻሕ ኣይተኻእለን”} ዝብል ኮይኑ፡ እዚ ማለት
  ላባራቶሪና ብዛዕባ ውጽኢት መርመራኹም ርግጸኛ ኣይኮነን ማለት’ዩ። ፍጹም ዝብሃል መርመራ የለን፡ ናብ መደብደምታ
  ክብጻሕ ዘይተኻእለ ውጽኢት ዘጋጥም ከኣ ንእሽቶ መጠን ናይ SARS-CoV-2 ኣብ ንረኽበሉ፡ ግን ከኣ ኮቪድ-19
  ኣለኩም ከብለና ዘይክእል መጠን ኣብ ንረኽበሉ እዋን እዩ።

\item

  \textbf{እንተድኣ መርመራኹም “ክምርመር ኣይተኻእለን”} ዝብል ኮይኑ፡ እዚ ማለት ኣብ መርኣያኹም
  ንሕና ክንምርምሮ እንኽእል እኹል ነገር ኣይነበረን፡ ወይ ነቲ መርመራ ኣብ ዘካየድናሉ እዋን ጌጋ ነይሩ ማለት’ዩ።
  ንመርኣያኹም ክንምርምሮ ስለዘይክኣልና፡ እቲ ኮቪድ-19 ዘስዕብ ቫይረስ ኣለኩም ድዩ የብልኩምን ክንረክብ ኣይክኣልናን።

\end{itemize}

\bigskip

\large \underline{\textbf{ምንጭታት}}

እንተድኣ ብዛዕባ ኮቪድ-19 ተወሳኺ ሓበሬታ ደሊኹም፡ ብኽብረትኩም ነዞም ካብ መሻርኽቲ ጥዕና ሕብረተሰብ ዝመጹ
ምንጭታት ርኣይዎም:

\begin{itemize}
\item
  Public Health \textemdash Seattle \& King County (ህዝባዊ ሓለዋ ጥዕና ሲያትልን ኪንግ
  ካውንትን):\\
  \url{https://www.kingcounty.gov/depts/health/covid-19/languages/tigrinya.aspx}

\item

  Washington State Department of Health (ናይ ክፍለ-ግዝኣት ዋሽንግተን ክፍሊ ጥዕና):\\
  \url{http://www.doh.wa.gov/Emergencies/Coronavirus}

\item
  U.S. Centers for Disease Control and Prevention (CDC, ማእከላት ቊጽጽርን ምግታእን ጥዕና
  መንግስቲ ኣመሪካ):\\
  \url{https://wwwn.cdc.gov/pubs/other-languages/?Sort=Lang%3A%3Aasc&Language=Tigrinya}

\end{itemize}

እንተድኣ ኣብ ተሌፎን ትዛረብዎ ሰብ ደሊኹም፡ ብኽብረትኩም ናብ ሓደ ካብዞም ዝስዕቡ ደውሉ:

\begin{itemize}

\item

  \textbf{ምስ ኮቪድ-19 ዝዛመዱ ሕክምናዊ ሕቶታት:} ነቲ King County Call Center (ማእከል
  ጻውዒት ተሌፎን ኪንግ ካውንቲ) ኣብ መንጎ ሰዓት 8 AM ከምኡ ድማ 7 PM በዚ ዝስዕብ ቁጽሪ ርኸብዎም
  \textbf{206-477-3977}

\item

  \textbf{ዘይሕክምናዊ ሕቶታት ብዛዕባ ኮቪድ-19፡ ንምስ ትካል-ዝዛመዱ ጒዳያት ዘጠቓለለ:} ነቲ King
  County COVID-19 Business and Community Information Line (ናይ ኪንግ ካውንቲ ናይ ትካላትን
  ማሕበረሰብን ናይ ሓበሬታ መስመር ተሌፎን)፡ ሰኑይ – ዓርቢ ኣብ መንጎ 8:30 AM ከምኡ ድማ 4:30 PM በዚ
  ዝስዕብ ቁጽሪ ርኸብዎ \textbf{206-296-1608}

\item

  \textbf{ሓፈሻዊ ሕቶታት ብዛዕባ ኮቪድ-19 ኣብ ክፍለ-ግዝኣት ዋሽንግተን:} ነቲ Washington State
  Novel Coronavirus Call Center (ማእከል ጻውዒት ተሌፎን ሓድሽ ኮረናቫይረስ ናይ ክፍለ-ግዝኣት
  ዋሽንግተን) ኣብ መንጎ 6 AM ከምኡ ድማ 10 PM በዚ ዝስዕብ ቁጽሪ ርኸብዎ \textbf{1-800-525-0127}

\end{itemize}

\bigskip

\large \underline{\textbf{ኣገባባትን ድሩትነታትን}}

\textbf{ኣገባባት}\\
እዚ መርመራ'ዚ ን SARs-CoV-2 ነቲ ግልባጥ-ቅዳሕ polymerase chain reaction (PCR, ፖሊመረይስ ናይ
ሰንሰለታዊ መልሰ-ተግባር) ተጠቒምካ ንምልላይ ዝተዳለወ እዩ። እዚ መርመራ'ዚ ምጒላሕ ናይቲ ቫይረስ ክፋላት ORf1B
ከምኡ ድማ S ይጥቀም ከምኡ'ውን ናይ ሰብ ቊጽጽር ምጒላሕ RnaseP ይጥቀም። ነቲ መርመራ ዘማዕበለን ናይ
ኣሰራርሕኡ ጠባያት ዝውስንን እቲ University of Washington’s Northwest Genomics Center እዩ።
እቲ ላባራቶሪ ብ WA Department of Health Medical Testing Site (ናይ ክፍሊ ጥዕና ዋሽንግተን
ማእከል ሕክምናዊ መርመራ) ልዑል ሕልኽላኽ ዘለዎ ክሊኒካዊ መርመራ ላባራቶሪ ንኸካይድ ፍቓድ ስራሕ ዝተዋህቦ እዩ።

\textbf{ድሩትነታት}

\begin{itemize}

\item

  እዞም ውጽኢታት እዚኦም ከም ኣካል ናይ ምርምር መጽናዕቲ ዝቐርቡ እዮም።

\item
  %- if staff_observed == "f"
  ነዚ መርኣያ'ዚ ወሲድኩምን ብገዛእ ርእስኹም (ወይ ንኻልእ ሰብ) ምልክት ገይርክምሉን ከምኡ ድማ ናብቲ ላባራቶሪ
  ብመገዲ ናይ ምብጽጻሕ ኣገልግሎት መሊስክሞ። ብበዓል ሞያ ጥዕና ከም ነርስ ወይ ዶክተር ዝተወስደ፡ ምልክት
  ዝተገብረሉ፡ ወይ ዝተገዓዓዘ ኣይኮነን።

  %- elif staff_observed == "t"
  ነዚ መርኣያ'ዚ ባዕልኻ (ወይ ንኻልእ ሰብ) ወሲድካዮ ምልክት ተገይርሉን ናብቲ ላባራቶሪ ከኣ ብመገዲ መልእኽተኛ
  ወይ ካልእ ዓይነት በዓል ሞያ ጥዕና ዘይኮነ ሰብ ተመሊሱን። ብከም ዶክተር ወይ ነርስ ዝኣመሰለ ወረቐት ስራሕ ዘለዎ
  በዓል ሞያ ኣይተወስደን ወይ ኣይተገዓዓዘን። 

  %- else
  \errmessage{Unknown value: staff_observed = "\VAR{staff_observed}"}

  %- endif

  እዚ ንውጽኢት መርመራኹም ክጸሉ ይኽእል፡ ወሃቢ ሕክምናዊ ኣገልግሎት ደጊሙ ክምርምረኩም ክደሊ ከኣ ይኽእል።

\item

  እዚ ጸብጻብ'ዚ ካብ ምርምር መጽናዕቲ ኣብ ዝተረኽበ ዳታ ዝምርኮስ ስለዝኾነ፡ ውጽኢታትኩም ህላውነት ኮቪድ-19
  ዘነጽሩ ዘይክኾኑ ይኽእሉ፡ እቲ ብረኽሲ ናይ SARS-CoV-2 ዝስዕብ ሕማም።

\end{itemize}

\end{document}
